\section{Related Work}\label{sec:related_work}
In this section, review the existing literature related to your study. Group the literature into logical categories/sub-topics.

\subsection{Sub-Topic 1 (e.g., Traditional Machine Learning Approaches)}
Discuss early, classic, or baseline approaches. Explain their strengths and point out their main limitations in your specific setting. Cite relevant works using author-year style: e.g., standard citation like \citep{lecun2015deep} or inline citations like \citet{vaswani2017attention}.

\subsection{Sub-Topic 2 (e.g., Deep Learning and Attention Models)}
Discuss more recent deep learning models, transformer-based approaches, or standard architectures related to your methodology. 

\subsection{Sub-Topic 3 (e.g., Domain-Specific Applications)}
If your work applies ML to a specific domain (e.g., computer vision, natural language processing, healthcare, finance), review the domain-specific models and explain how they relate to your study.

\begin{table}[h]
\caption{Comparison of our proposed method with existing approaches. ($\checkmark$ indicates the feature is supported, while $\times$ indicates it is not.)}\label{tab:related_work_comparison}
\begin{tabular}{@{}lcccc@{}}
\toprule
Method & Scalability & Robustness & Interpretability & Computational Cost \\
\midrule
Method A \citep{lecun2015deep} & $\checkmark$ & $\times$ & $\times$ & High \\
Method B \citep{vaswani2017attention} & $\times$ & $\checkmark$ & $\times$ & Very High \\
Method C \citep{goodfellow2016deep} & $\checkmark$ & $\checkmark$ & $\times$ & High \\
\textbf{Ours (Method Name)} & \textbf{$\checkmark$} & \textbf{$\checkmark$} & \textbf{$\checkmark$} & \textbf{Low} \\
\bottomrule
\end{tabular}
\end{table}
